% !TEX program = xelatex
% Modern Blue One-page Resume Template
% Compile with XeLaTeX on Overleaf.
% Put your portrait image as photo.jpg in the same directory if you want a photo.
% If photo.jpg is missing, a placeholder box will be shown automatically.

\documentclass[11pt,a4paper]{article}

% ==================== Font / Layout ====================
\usepackage[UTF8,fontset=fandol]{ctex}
\usepackage{fontspec}

% Font fallback:
% English: Inter -> IBM Plex Sans -> default
% Chinese: Noto Sans CJK SC -> Source Han Sans SC -> Fandol
\IfFontExistsTF{Inter}{
  \setmainfont{Inter}
}{
  \IfFontExistsTF{IBM Plex Sans}{
    \setmainfont{IBM Plex Sans}
  }{}
}

\IfFontExistsTF{Noto Sans CJK SC}{
  \setCJKmainfont{Noto Sans CJK SC}
}{
  \IfFontExistsTF{Source Han Sans SC}{
    \setCJKmainfont{Source Han Sans SC}
  }{}
}

\usepackage[left=0.90cm,right=0.90cm,top=0.56cm,bottom=0.56cm]{geometry}
\usepackage{tabularx}
\usepackage{array}
\usepackage{enumitem}
\usepackage{titlesec}
\usepackage{xcolor}
\usepackage{hyperref}
\usepackage{bookmark}
\usepackage{fontawesome5}
\usepackage{graphicx}
\usepackage[most]{tcolorbox}

% ==================== Color System ====================
\definecolor{accent}{HTML}{2563EB}
\definecolor{body}{HTML}{1F2937}
\definecolor{muted}{HTML}{6B7280}

% ==================== Hyperlink / PDF Metadata ====================
% Change metadata before publishing.
\hypersetup{
  colorlinks=true,
  urlcolor=accent,
  linkcolor=accent,
  pdfborder={0 0 0},
  pdftitle={YOUR NAME Resume},
  pdfauthor={YOUR NAME},
  pdfsubject={Resume},
  pdfkeywords={
    Resume,
    AI Agent,
    LLM,
    RAG,
    Deep Learning,
    Machine Learning,
    Python,
    PyTorch
  }
}

% ==================== Global Style ====================
\pagestyle{empty}
\setlength{\parindent}{0pt}
\setlength{\parskip}{0pt}
\renewcommand{\arraystretch}{1.06}

\setlist[itemize]{
  leftmargin=1.15em,
  itemsep=0.75pt,
  topsep=1pt,
  parsep=0pt,
  partopsep=0pt,
  label={\raisebox{0.2ex}{\scriptsize\textcolor{accent}{$\bullet$}}}
}

% ==================== Section Style ====================
\titleformat{\section}
  {\large\bfseries\color{accent}}
  {}{0em}
  {\textcolor{accent}{\rule[-1pt]{3pt}{13pt}}\hspace{6pt}}
  [{\vspace{-2pt}\color{accent!45}\titlerule[0.45pt]}]

\titlespacing*{\section}{0pt}{5pt}{2.5pt}

% ==================== Commands ====================
\newcommand{\cvdate}[1]{\hfill {\small\color{muted}#1}}
\newcommand{\cvmeta}[1]{{\small\color{muted}#1}}
\newcommand{\cvplaceholder}[1]{{\color{muted}#1}}

% ATS-friendly section:
% #1: visible Chinese/English title
% #2: PDF outline title for parsers
% #3: stable anchor name, no spaces
\newcommand{\cvsection}[3]{%
  \phantomsection
  \pdfbookmark[1]{#2}{sec:#3}
  \section*{#1}
}

% Optional project GitHub link.
% Leave the 4th argument empty if no repo link is available.
\newcommand{\projectgit}[1]{%
  \if\relax\detokenize{#1}\relax
  \else
    \hspace{4pt}\href{#1}{\textcolor{accent}{\footnotesize\faGithub}}%
  \fi
}

% Education entry:
% \cvedu{School}{Degree / Major}{Time}
\newcommand{\cvedu}[3]{%
  \textbf{#1} \cvdate{#3}\\[-1pt]
  \cvmeta{#2}\vspace{1pt}
}

% Project entry:
% \cvproject{Project Name}{Time}{Tech Stack}{GitHub URL or empty}
\newcommand{\cvproject}[4]{%
  \vspace{1.5pt}
  \textbf{#1}\projectgit{#4} \cvdate{#2}\\[-1pt]
  \cvmeta{\textcolor{accent}{\rule[0.45ex]{2.4pt}{0.8ex}}\hspace{5pt}#3}
}

% Research / open-source item:
% \cvresearch{Title}{Time}{Subtitle / Context}
\newcommand{\cvresearch}[3]{%
  \vspace{1.5pt}
  \textbf{#1} \cvdate{#2}\\[-1pt]
  \cvmeta{\textcolor{accent}{\rule[0.45ex]{2.4pt}{0.8ex}}\hspace{5pt}#3}
}

% Photo. Upload photo.jpg to project root.
\newcommand{\cvphoto}{%
  \IfFileExists{photo.jpg}{%
    \fcolorbox{accent!45}{white}{%
      \includegraphics[width=2.15cm,height=2.72cm,keepaspectratio]{photo.jpg}%
    }%
  }{%
    \fcolorbox{accent!45}{white}{%
      \parbox[c][2.72cm][c]{2.12cm}{%
        \centering\small\color{muted}Photo%
      }%
    }%
  }%
}

% Skill tag pill.
\newtcbox{\skilltag}{
  on line,
  box align=base,
  colback=accent!7,
  colframe=accent!18,
  coltext=accent,
  fontupper=\scriptsize,
  arc=2.2pt,
  boxrule=0.25pt,
  left=3pt,
  right=3pt,
  top=1pt,
  bottom=1pt,
  nobeforeafter
}

% Skill row:
% \skillrow{Category}{\skilltag{A} \skilltag{B}}
\newcommand{\skillrow}[2]{%
  \textbf{\color{accent}#1}\quad #2\par\vspace{1.1pt}
}

% ==================== Document ====================
\begin{document}
\color{body}
\fontsize{9.95pt}{11.35pt}\selectfont

% ATS-friendly top-level PDF bookmark.
\phantomsection
\pdfbookmark[0]{Resume - YOUR NAME}{resume}

% ==================== Header ====================
\begin{tabularx}{\linewidth}{@{}X m{2.42cm}@{}}
  \begin{minipage}[c]{\linewidth}
    {\fontsize{28pt}{30pt}\selectfont\bfseries\color{accent} YOUR NAME}\\[6pt]
    {\Large\bfseries Target Role / Target Internship}\\[5pt]

    {\small
    \begin{tabularx}{\linewidth}{@{}>{\bfseries}l>{\raggedright\arraybackslash}X>{\bfseries}l>{\raggedright\arraybackslash}X@{}}
    Email: & \href{mailto:your.email@example.com}{your.email@example.com}
    & Phone: & (+86) XXX-XXXX-XXXX \\[2pt]

    GitHub: & \href{https://github.com/yourname}{github.com/yourname}
    & Location: & City / Remote \\[2pt]

    Target City: & City A / City B / Remote
    & Availability: & Immediately / 5 days per week \\
    \end{tabularx}
    }
  \end{minipage}
  &
  \begin{center}
    \cvphoto
  \end{center}
\end{tabularx}

\vspace{-5pt}

% ==================== Education ====================
\cvsection{Education}{Education}{education}

\cvedu
{University Name}
{Master of Science, Major / Track}
{2026.01 -- 2028.02 Expected}

\cvmeta{Core Courses: Machine Learning, Deep Learning, Optimization, Data Mining}

\vspace{2pt}

\cvedu
{University Name}
{Bachelor of Engineering, Computer Science and Technology}
{2020.09 -- 2024.06}

\cvmeta{Core Courses: Data Structures, Computer Networks, Computer Organization, Operating Systems, Discrete Mathematics, Linear Algebra}

\vspace{2pt}

% ==================== Projects ====================
\cvsection{Projects}{Projects}{projects}

\cvproject
{Project One: Short Descriptive Title}
{2025.12 -- 2026.03}
{Python · PyTorch · LangGraph · ChromaDB · Docker}
{https://github.com/yourname/project-one}
\begin{itemize}
  \item Describe the project goal, scenario, and technical architecture in one sentence.
  \item Explain your implementation details and emphasize frameworks, models, algorithms, or engineering decisions.
  \item Add measurable results if available, such as latency, accuracy, recall, GPU usage, dataset size, or cost reduction.
\end{itemize}

\cvproject
{Project Two: Short Descriptive Title}
{2026.03 -- 2026.04}
{Python · FastAPI · SQLite · MCP SDK · API Integration}
{}
\begin{itemize}
  \item Describe what problem the project solves and what users or systems it serves.
  \item Highlight the core components you built, such as tools, resources, prompts, services, or pipelines.
\end{itemize}

\cvproject
{Project Three: Short Descriptive Title}
{2026.04 -- 2026.05}
{Python · LLM API · ReAct · RAG · External APIs}
{}
\begin{itemize}
  \item Describe the agent or application workflow: planning, retrieval, reasoning, tool use, and output generation.
  \item Explain how you handled reliability, observability, fallback, cost control, or multi-provider integration.
\end{itemize}

\cvproject
{Other Project: Compact Description}
{2025}
{Frontend · Backend · Database · LLM API}
{}
\begin{itemize}
  \item Keep secondary projects short; one bullet is usually enough for a one-page resume.
\end{itemize}

% ==================== Research / Open Source ====================
\cvsection{Research \& Open Source}{Research and Open Source}{research-open-source}

\begin{itemize}
  \item \textbf{Research / Implementation:} Summarize a paper reproduction, algorithm study, or course research project. Add dataset, model, metric, and result where possible.
  \item \textbf{Open Source:} Mention a template, tool, library, or repository you built. Add a GitHub link if available.
  \item \textbf{Competition / Award:} Add ICPC / Kaggle / school award / scholarship, if relevant. Keep it concise.
\end{itemize}

% ==================== Skills ====================
\cvsection{Skills}{Skills}{skills}

\skillrow{Programming}{
  \skilltag{Python}
  \skilltag{C/C++}
  \skilltag{JavaScript / TypeScript}
  \skilltag{SQL}
}

\skillrow{AI / ML}{
  \skilltag{PyTorch}
  \skilltag{Transformer}
  \skilltag{RAG}
  \skilltag{LLM API}
  \skilltag{Agent}
  \skilltag{Diffusion Models}
}

\skillrow{Engineering}{
  \skilltag{Linux / macOS}
  \skilltag{Git}
  \skilltag{Docker}
  \skilltag{FastAPI}
  \skilltag{Gradio}
  \skilltag{LaTeX}
}

\skillrow{Languages}{
  \skilltag{Chinese Native}
  \skilltag{English IELTS / CET}
  \skilltag{Technical Reading and Writing}
}

\end{document}
